\documentclass[10pt,a4paper]{article}

\usepackage[utf8]{inputenc}
\usepackage[T1]{fontenc}
\usepackage[spanish]{babel}
\usepackage[margin=2cm]{geometry}
\usepackage{graphicx}
\usepackage{amsmath, amssymb}
\usepackage{subcaption}
\usepackage{float}
\usepackage[hidelinks]{hyperref}
\usepackage{booktabs}
\usepackage{algorithm}
\usepackage{algpseudocode}
\usepackage{caption}
\usepackage{xcolor}
\usepackage{multicol}
\usepackage{tabularx}
\setlength{\tabcolsep}{3pt}

\newcommand{\twopaneltitles}[2]{%
    \noindent
    \begin{minipage}[t]{0.49\linewidth}
        \centering\normalsize\bfseries #1
    \end{minipage}\hfill%
    \begin{minipage}[t]{0.49\linewidth}
        \centering\normalsize\bfseries #2
    \end{minipage}\par\vspace{2pt}%
}

\begin{document}

\begingroup
\centering
{\footnotesize\scshape\color{black!70} Universidad de San Andrés \quad$\cdot$\quad Escuela de Ingeniería \quad$\cdot$\quad I308 --- Visión Artificial}\\[4pt]
{\color{black!25}\rule{\linewidth}{0.6pt}}\\[10pt]
{\fontsize{20}{24}\selectfont\bfseries\color{black}
Reconstrucción 3D\\[2pt]
por estéreo multivista\par}
\vspace{6pt}
{\normalsize\itshape\color{black!70} Trabajo Práctico II\par}
\vspace{10pt}
{\normalsize\bfseries Agustín Patruno Carranza \quad$\cdot$\quad Catalina Hirsch \quad$\cdot$\quad Mathias Kneipp\par}
\vspace{2pt}
{\footnotesize\color{black!60}\today\par}
\vspace{4pt}
{\color{black!25}\rule{\linewidth}{0.6pt}}
\endgroup
\vspace{10pt}

\begin{abstract}
\noindent

\end{abstract}

\begin{multicols}{2}

\section{Introducción y Marco Teórico}
% Introducir los conceptos fundamentales de forma clara y concisa.

\subsection{Modelo de cámara pinhole y distorsión}
% ¿Cuáles son los parámetros intrínsecos? ¿Por qué es necesario calibrar?

\subsection{Calibración estéreo}
% ¿Qué parámetros extrínsecos relativos se estiman y para qué sirven?
% ¿Cómo se realiza esta estimación? ¿Qué es el error de reproyección?

\subsection{Geometría epipolar y rectificación}
% ¿Qué problema resuelve la rectificación? ¿Cómo facilita el cálculo de disparidad?

\subsection{De disparidad a profundidad}
% ¿Qué relación matemática hay entre disparidad y profundidad?

\section{Metodología}
% Describir brevemente el pipeline implementado, mencionando en cada etapa
% las técnicas y parámetros utilizados.

\subsection{Calibración y rectificación}
% Herramienta utilizada (tool de la cátedra o funciones de OpenCV).

\subsection{Disparidad y nube local}
% Método elegido (SGBM, red neuronal, etc.) y parámetros ajustados.
% Cómputo de la nube local con reprojectImageTo3D (matriz Q).

\subsection{Pose global y transformación}
% ¿Cómo estimaron la pose con solvePnP? ¿Qué representa el vector de Rodrigues?
% ¿Qué cuenta hicieron para llevar la nube de puntos al sistema global?

\subsection{Fusión y optimización en Open3D}
% Pasos de optimización usados y parámetros elegidos justificados
% (bounding box -> voxel_down_sample -> remoción de outliers).

\subsection{Medición métrica}
% ¿Cómo midieron el objeto? ¿Por qué es posible tener un resultado métrico
% (y no el problema de reconstrucción salvo por una escala)?

\section{Resultados y Discusión}

\subsection{Resultados}
% Evidencia visual y cuantitativa de cada etapa del pipeline:
% - Calibración: parámetros obtenidos, error RMS, par original vs. rectificado.
% - Disparidad: mapa de disparidad de un par representativo, nube local.
% - Pose y marco global: poses de cámara superpuestas con el marco de referencia.
% - Nube global y optimización: antes/después en cada paso, cantidad de puntos.
% - Medición métrica: tamaño estimado en mm, método, error relativo si se conoce.

\subsection{Discusión}
% Análisis crítico: problemas encontrados en el pipeline, limitaciones del método
% y posibles mejoras.

\begin{thebibliography}{9}
\footnotesize
\bibitem{hartley} R. Hartley y A. Zisserman. \href{https://www.robots.ox.ac.uk/~vgg/hzbook/index.html}{Multiple View Geometry in Computer Vision}. 2.ª ed., Cambridge University Press, 2004.
\bibitem{opencv} OpenCV. \href{https://docs.opencv.org/4.x/d9/d0c/group__calib3d.html}{Camera Calibration and 3D Reconstruction}. Documentación oficial.
\bibitem{open3d} Q.-Y. Zhou, J. Park y V. Koltun. \href{http://www.open3d.org/}{Open3D: A Modern Library for 3D Data Processing}. arXiv:1801.09847, 2018.
\end{thebibliography}

\end{multicols}
\end{document}
